\clearpage
\hypertarget{native-470-report}{}
\begin{center}
\rule{0.55\linewidth}{0.4pt}\\[2pt]
Astronomische Nachrichten. Vol. XXI. No. 496. p. 247. 17 February 1844.\\[2pt]
\rule{0.55\linewidth}{0.4pt}
\end{center}

In a letter of the 16th, Councillor \textsc{Gauss} sent me the following observations of the comet:
\begin{center}
\setlength{\tabcolsep}{3pt}
\begin{tabular}{l|l|l|l}
1844 January 15 & $8^{\rm u}18^{\rm m}17^{\rm s}8$ mean time & R.A.: $77^\circ20'59''6$ & $\delta$: ---\quad ---\quad --- \\
 & $8\;10\;44.5$\quad --- & R.A.: ---\quad ---\quad --- & $\delta$: $+3^\circ49'47''3$ \\
\end{tabular}
\end{center}

For that day, Dr. \textsc{Goldschmidt} had calculated a star’s position as $78^\circ17'2''3$ $+3^\circ51'2''9$. At $8^{\rm u}18^{\rm m}17^{\rm s}8$ the comet preceded that star by $3'34''18$, and at $8^{\rm u}10^{\rm m}44^{\rm s}5$ mean time it was $1'15''6$ farther south. Councillor \textsc{Gauss} considers this observation as good as he could make with the ring micrometer given how faint the comet was, provided, as goes without saying, that the star’s position is well determined.

In a later letter of the 21st, Councillor \textsc{Gauss} kindly sent me the comparison with \textsc{Goldschmidt}’s second set of elements.
\begin{flushright}\textsc{Schumacher.}\end{flushright}
